\sezione{12}{Quando usarlo, quando no}{Parte III · pregi, limiti e checklist finale}

\begin{inbreve}
Il Random Forest è un ottimo modello «di partenza» per dati tabellari: robusto, con pochi parametri, poco sensibile alla scala delle variabili e con una stima dell'errore gratuita (OOB). Ha però limiti chiari: modelli pesanti, interpretazione meno diretta di un singolo albero, nessuna estrapolazione.
\end{inbreve}

\subsection{Pregi e limiti}

\begin{center}
\begin{tcbraster}[raster columns=2,raster equal height,raster column skip=10pt]
\begin{tcolorbox}[enhanced,unbreakable,arc=4pt,boxrule=0pt,colback=verde!7,borderline west={3pt}{0pt}{verde},
  title={\ding{51}\ \ Pregi},fonttitle=\headfont\bfseries\small,colbacktitle=verde,coltitle=white,fontupper=\small]
\begin{itemize}
  \item Ottime prestazioni con i default.
  \item Gestisce interazioni e non linearità senza doverle specificare.
  \item Non serve standardizzare le feature.
  \item Robusto a feature irrilevanti e a outlier nelle feature.
  \item Stima OOB senza validation set.
  \item Addestramento parallelo (\texttt{n\_jobs=-1}).
  \item Classificazione e regressione con la stessa API.
\end{itemize}
\end{tcolorbox}
\begin{tcolorbox}[enhanced,unbreakable,arc=4pt,boxrule=0pt,colback=rosso!6,borderline west={3pt}{0pt}{rosso},
  title={\ding{55}\ \ Limiti},fonttitle=\headfont\bfseries\small,colbacktitle=rosso,coltitle=white,fontupper=\small]
\begin{itemize}
  \item Modelli grandi in memoria e previsione più lenta di un modello lineare.
  \item Meno interpretabile di un singolo albero o di una regressione.
  \item In regressione non estrapola oltre i target visti.
  \item Le importanze MDI sono distorte; con feature correlate serve cautela.
  \item Su dati tabellari il gradient boosting ben regolato spesso è un po' più accurato.
  \item Poco adatto a immagini, audio, testo grezzo.
\end{itemize}
\end{tcolorbox}
\end{tcbraster}
\end{center}

\subsection{Random Forest e Gradient Boosting}

\begin{mlbox}[Due modi diversi di usare tanti alberi]
\begin{itemize}
  \item \hl{Random Forest} (\emph{bagging}): alberi \key{profondi e indipendenti}, addestrati in parallelo, poi messi in media. Riduce la \emph{varianza}. Difficile da «rompere» con iperparametri sbagliati.
  \item \hl{Gradient Boosting} (in scikit-learn \texttt{HistGradientBoostingClassifier/Regressor}): alberi \key{piccoli e sequenziali}, ognuno corregge gli errori dei precedenti. Riduce soprattutto il \emph{bias}. Spesso più accurato, ma più sensibile al tuning (learning rate, numero di iterazioni).
\end{itemize}
Una strategia pratica: prova prima un modello lineare, poi un Random Forest con i default, poi un gradient boosting; confrontali con la stessa validazione incrociata.
\end{mlbox}

\subsection{Problemi frequenti e soluzioni}

\begin{center}
\small
\renewcommand{\arraystretch}{1.4}
\rowcolors{2}{chiaro}{white}
\begin{tabularx}{\linewidth}{>{\raggedright\arraybackslash}p{4.4cm} >{\raggedright\arraybackslash}p{4.2cm} >{\raggedright\arraybackslash}X}
\rowcolor{navy}\thd{Sintomo} & \thd{Causa probabile} & \thd{Soluzione}\\
Accuratezza di training 100\%, test molto più bassa & Normale per alberi pieni; se il test è basso, rumore o pochi dati & Guarda OOB/CV, non il training; aumenta \texttt{min\_samples\_leaf}\\
Classe rara mai prevista & Classi sbilanciate & \texttt{class\_weight="balanced"}, soglia su \texttt{predict\_proba}, metriche come recall/AUC\\
File del modello enorme & Alberi pieni su molti dati & \texttt{max\_depth}, \texttt{min\_samples\_leaf}, meno alberi\\
\texttt{ValueError} con colonne di testo & Feature non numeriche & \texttt{OneHotEncoder} in un \texttt{ColumnTransformer}\\
Risultati diversi a ogni esecuzione & Seme non fissato & \texttt{random\_state} nel modello, nello split e nella CV\\
Prestazione ottima in CV, scarsa in produzione & Data leakage o dati cambiati & Preprocessing nella pipeline; controlla le feature disponibili al momento della previsione\\
Warning sull'OOB / \texttt{NaN} nell'OOB & Troppi pochi alberi & Aumenta \texttt{n\_estimators}\\
\end{tabularx}
\end{center}

\subsection{Checklist finale}

\begin{passi}
  \item Ho separato il test set con \texttt{stratify} (classificazione) e fissato \texttt{random\_state}?
  \item Il preprocessing (categorie, mancanti) è dentro una \texttt{Pipeline}?
  \item Ho confrontato il Random Forest con un modello semplice usando la stessa CV?
  \item Ho scelto \texttt{max\_features} e la dimensione delle foglie con la CV (o almeno con l'OOB)?
  \item Ho usato metriche adatte al problema (recall, AUC, MAE\dots) e non solo l'accuratezza?
  \item Per interpretare il modello ho usato l'importanza per permutazione su dati non visti?
  \item Ho salvato la pipeline intera con la versione di scikit-learn usata?
\end{passi}

\begin{appunti}{Riepilogo del capitolo}
\begin{itemize}
  \item Usalo come modello robusto di partenza su dati tabellari.
  \item Limiti: memoria, interpretabilità, nessuna estrapolazione, importanze MDI distorte.
  \item Bagging (RF) riduce la varianza; boosting riduce il bias: confrontali.
  \item Segui la checklist a ogni progetto.
\end{itemize}
\end{appunti}
